\section{AI Design}
\label{sec:ai}

Our AI system employs different algorithms for each phase, reflecting the
fundamentally different information structures.

\subsection{Phase 2: Alpha-Beta Minimax}

Phase~2 is a perfect-information game: through card counting during Phase~1,
both players' hands are fully known. We solve Phase~2 positions exactly using
alpha-beta pruning with the following optimizations:

\begin{itemize}
  \item \textbf{Bitmask representation:} Each hand is encoded as a 64-bit
    integer, where bit $i$ represents card $i$ (using the mapping
    $\text{id} = (\text{rank} - 2) \times 4 + \text{suit}$). This enables
    $O(1)$ card addition/removal via XOR.
  \item \textbf{Move ordering:} Leaders play high cards first (descending
    rank); followers try cheap wins first (ascending rank). Good move ordering
    is critical for alpha-beta efficiency.
  \item \textbf{Transposition table:} Positions are cached by
    $(\text{mask}_0, \text{mask}_1, \text{leader}, \text{lead\_id})$ to avoid
    re-solving identical positions reached via different move orders.
  \item \textbf{C++ implementation:} The solver is implemented in C++ via
    pybind11, using Zobrist hashing for the transposition table and iterative
    deepening for move ordering. All 13-card positions solve exactly in
    under 75ms.
\end{itemize}

\begin{algorithm}[H]
\caption{Alpha-Beta Minimax for Phase 2}
\KwIn{hands $H_0, H_1$; trump $\tau$; leader $\ell$; bounds $\alpha, \beta$}
\KwOut{(score for Player~0, best card)}
\If{both hands empty}{
  \Return $(0, \text{null})$
}
\If{in transposition table}{
  \Return cached value
}
\ForEach{legal card $c$ in ordered moves}{
  Play $c$, resolve trick if both cards played\;
  $\text{score} \leftarrow \text{Minimax}(\text{new state}, \alpha, \beta)$\;
  Update best move and $\alpha/\beta$ bounds\;
  \If{$\beta \leq \alpha$}{
    \textbf{break} \tcp{Prune remaining branches}
  }
}
Store in transposition table\;
\Return (best score, best card)
\end{algorithm}

\subsubsection{Quick-Trick Evaluation}

For positions too large for exact solving, we estimate Player~0's tricks using
a \emph{quick-trick} heuristic that counts:
\begin{enumerate}
  \item \textbf{Top winners:} Consecutive top cards in each suit (e.g., if we
    hold the Ace and King of a suit and no higher cards remain, both are
    winners).
  \item \textbf{Long-suit promotions:} Extra cards in suits where we have
    more than the opponent, which become winners after the opponent's cards
    are exhausted.
  \item \textbf{Ruffing potential:} Void suits with trump holdings enable
    winning tricks by trumping opponent leads.
  \item \textbf{Lead advantage:} Having winners when on lead is more valuable
    than off-lead winners.
\end{enumerate}

\subsection{Phase 1: Rule-Based Heuristic}

Phase~1 presents an imperfect-information challenge: we don't know the
opponent's hand or the stock order. We use a rule-based heuristic that
decomposes each trick into two decisions:

\begin{enumerate}
  \item \textbf{Win or lose?} The face-up card is evaluated against a
    threshold: always win Aces and Kings, win trump cards ranked 7 or higher,
    win Queens and Jacks of any suit, and lose deliberately for low cards.
  \item \textbf{Card selection:} When trying to win, play the cheapest card
    that beats the opponent. When trying to lose, shed the lowest non-trump
    card. When leading to win, lead from the strongest suit; to lose, lead
    the lowest non-trump.
\end{enumerate}

This simple strategy proves highly effective: it wins 93\% of games against
random play and provides a strong hand for Phase~2.

\subsubsection{Card Counting}

Throughout the game, the AI maintains a \emph{card counter} that tracks:
\begin{itemize}
  \item All cards in its own hand
  \item All cards played in tricks
  \item All face-up cards observed (including those taken by the opponent)
  \item The set of cards whose location is unknown
\end{itemize}

At the transition to Phase~2, the card counter deduces the opponent's exact
hand: it is precisely the set of all 52 cards minus those in our hand and
those that have been played. This enables exact minimax in Phase~2.

\subsection{Difficulty Levels}

\begin{center}
\begin{tabular}{lll}
\toprule
Level & Phase 1 & Phase 2 \\
\midrule
Easy & Rule-based heuristic & Rule-based heuristic \\
Medium/Hard & Rule-based heuristic + card counting & Alpha-beta minimax \\
\bottomrule
\end{tabular}
\end{center}

The primary difference between Easy and Medium/Hard is the Phase~2 solver:
the heuristic uses simple rules (lead aces, play long suits), while the AI
uses exact minimax for positions with $\leq 10$ cards per player.
